\documentclass[french]{llncs}

\usepackage{erewhon}
\usepackage[utf8]{inputenc}
\usepackage[french]{babel}
\usepackage{csquotes}
\usepackage{caption}

\addto\captionsfrench{%
  \renewcommand{\theoremname}{Théorème}%
}
\addto\captionsfrench{%
  \renewcommand{\lemmaname}{Lemme}%
}
\addto\captionsfrench{%
    \renewcommand{\keywordname}{\textbf{Mots-clés~:}}%
}

% Bibliography.
\usepackage[natbib=true,backend=bibtex]{biblatex}
\addbibresource{../refs.bib}

% Mathematics stuff.
\usepackage{amsmath}
\usepackage{amssymb}
\usepackage{xspace}
\usepackage{pgfplots}
\usepackage{pgfplotstable}
\usepgfplotslibrary{groupplots}
\pgfplotsset{compat=1.18}
\usepackage[ruled,vlined]{algorithm2e} % pseudo-code.

% TikZ
\usepackage{tikz}
\usepackage{pgfplots}

% Miscellaneous
\usepackage{multicol}
\usepackage{listings}
\usepackage{xcolor} % Driver-independent color extensions for LaTeX
\usepackage[inline,shortlabels]{enumitem}

\newcommand{\theauthor}{A. CN}
\newcommand{\thetitle}{Démonstration de la résilience d'un réseau P2P}

% Metadata
\usepackage{hyperref}
\hypersetup{
    pdfauthor={\theauthor},
    pdftitle={\thetitle},
}

\begin{document}

\title{\thetitle}
\author{\theauthor}
\authorrunning{A. CN}
\institute{Gravitalia\\
\email{acn@gravitalia.com}}

\maketitle

\keywords{Résilience \and Pair-à-pair \and Théorie des graphes \and Théorie de 
la percolation}

\begin{abstract}
    Cette démonstration, par la théorie des graphes, permet d'aborder la 
capacité des réseaux pair-à-pair à tolérer la disparition de plusieurs 
n\oe{}uds.
\end{abstract}

\section{Introduction}

La démonstration modélise un graphe aléatoire de sommets (n\oe{}uds) 
éventuellement connectés entre eux. Chaque sommet représente un pair. Cette 
démonstration modélise un réseau pair-à-pair, où chaque pair agit au bénéfice 
des autres~; et non un réseau point-à-point. Nous considérons qu'il n'existe ni 
congestion réseau ni attaque sur le réseau.

\section{Démonstration}

Soit une famille de variables de Bernoulli indépendantes 
\(\left(X_{i,j}\right)_{1 \le i < j \le n}\) de paramètre \(p \in 
\left[0,1\right]\), où \(n \in \mathbb{N}\) est le nombre de sommets, 
définissant un graphe d'Erdös-Rényi\cite{erdos1960} \(G(n,p)\) tel que 
\(X_{i,j}=1\) si l'arête \((i,j)\) existe et \(0\) sinon.

\nocite{gayral2017} \nocite{thevenin2016} % Thanks for their lessons!

\begin{theorem}[Connexité]\label{thm:connected}
  \(p(n) = \frac{\log n + c_n}{n}\) est seuil critique pour la connexité du 
graphe d'Erdős-Rényi.
  Plus précisément,
  \[
      \mathbb{P}(G(n,p(n)) \text{ est connexe}) 
\mathop{\to}\limits_{c_n\to\pm\infty}
      \begin{cases}
          1 \text{ si } c_n \to +\infty\\
          0 \text{ si } c_n \to -\infty
      \end{cases}
  \]
\end{theorem}

\begin{lemma}[Composante géante]\label{lem:geant}
  Posons \(p = \frac{c}{n}\), avec \(c = n \cdot p\) la densité moyenne de 
connexions par sommet.
  \begin{itemize}
      \item Si \(c<1\), alors en moyenne chaque sommet à moins d'un voisin. Les 
composantes sont de tailles \(\Theta(\log n)\).
      \item Si \(c=1\), alors il y a une transition du seuil de percoation. La 
plus grande composante est de taille \(\Theta(n^{2/3})\).
      \item Si \(c>1\), alors une composante géante apparaît et contient une 
fraction significative des n\oe{}uds.
  \end{itemize}
\end{lemma}

A partir du lemme \ref{lem:geant}, on peut déterminer lorsqu'un réseau 
pair-à-pair est viable, et s'assurer que les données arrivent au pair voulu. 
Seuls \(\alpha~\%\) des n\oe{}uds sont en dehors de la composante géante, avec :
\[ \alpha = e^{-c(1-\alpha)} \]

\begin{theorem}[de Menger]\label{thm:menger}
    Soient \(G\) un graphe non-orienté, et \(x\) et \(y\) deux sommets 
distincts.
    Le nombre maximal de chemins  \(x\mbox{-}y\) internement disjoints est égal 
au nombre minimal de sommets dont la suppression sépare \(x\) de \(y\).
\end{theorem}

Dans le contexte du réseau P2P, le nombre maximal de chemins disjoints entre le 
pair émetteur \(x\) et le pair récepteur \(y\) est égal à la résilience 
minimale de la connexion \((x,y)\) contre la défaillance des intermédiaires.

Le théorème \ref{thm:connected} nous donne le seuil pour la 1-connexité 
(connexité simple, où \(k=1\)) du graphe total.
Pour assurer une résilience, nous avons besoin d'une k-connexité avec 
\(k\geq2\). Le seuil de probabilité \(p(n)\) requis pour qu'un graphe 
d'Erdős-Rényi soit k-connexe est :

\begin{theorem}[Seuil de k-Connexité]\label{thm:k-connected}
    Pour tout entier fixe \(k\geq1\), la probabilité d'arête \(p(n)\) tel que
    \begin{equation*}
        p(n) \geq \dfrac{\log n + (k-1)\log(\log n)+\omega(1)}{n}
    \end{equation*}
    est un seuil critique pour la k-connexité de \(G(n,p(n))\).
\end{theorem}

Partons des postulats que les serveurs d'un réseau (\(m\)) centralisé soient 
entièrement connectés (\(p_m=1\)) et que \(m\) croît plus lentement que 
\(n\)\footnote{Par exemple, on peut avoir \(m = \Theta(\frac{\log n}{\log \log 
n})\). Empiriquement, on observe que \(m\) croît plus lentement que \(n\). Ce 
postulat permet de garantir une solution pour un \(n\) très grand.}. On cherche 
donc quand un réseau pair-à-pair est plus résilient qu'un réseau centralisé, 
sachant que \(n>m\), c'est-à-dire que dans les deux cas, il y a plus 
d'utilisateurs que de serveurs -- ce qui est le cas empiriquement.

Notons \(\mathcal{K}\) la fonction associé au théorème \ref{thm:menger}. Nous 
cherchons donc lorsque :

\[
    \mathcal{K}(G_{p2p}) \geq \mathcal{K}(G_{m})
\]

Ainsi, pour déconnecter deux n\oe{}uds dans \(\mathcal{K}_m\), il faut retirer 
les \(m-2\) autres n\oe{}uds, donc la connectivité de \(\mathcal{K}_m\) est :

\begin{equation}
    \label{eq:server}
    \mathcal{K}(G_{\text{serveur}}) = \mathcal{K}(G(m,1))=m-1
\end{equation}

Donc, pour que le réseau pair-à-pair soit au moins résilient que le réseau 
centralisé, il faut que :
\[
    \mathcal{K}(G_{p2p}) \geq m-1
\]

Injectons (\ref{eq:server}) dans la condition obtenue au théorème 
\ref{thm:k-connected}.

\[
    p(n)\gtrsim \dfrac{\log n + ((m(n)-1)-1)\log(\log n)+\omega(1)}{n} = 
        \boxed{\dfrac{\log n + (m(n)-2)\log(\log n)+\omega(1)}{n}}
\]

\paragraph{Conclusion.} Bien que le réseau pair-à-pair ait plus de n\oe{}uds, 
il n'est pas plus résilient par nature. Il ne le devient que lorsque sa densité 
de connexion \(p(n)\) est suffisamment élevée pour atteindre le seuil de 
\(m\)-connexité, garantissant au moins \(m\) chemins disjoints entre les pairs, 
surpassant ainsi les \(m-1\) chemins disjoints du réseau centralisé 
\(\mathcal{K}_{m}\).

\paragraph{Critique du modèle.} Nous nous sommes néanmoins intéressé à un cas 
très favorable au réseau centralisé, représenté sous forme d'un graphe complet. 
Dans la pratique, tous les serveurs ne sont pas connectés entre eux, donc 
\(p_m<1\). Le réseau pair-à-pair est donc résilient plus rapidement que le 
réseau centralisé. Néanmoins, nous omettons la différence entre suppression 
aléatoire et suppression stratégique. Dans ce dernier cas, les n\oe{}uds les 
plus connectés sont supprimés et la résilience du réseau pair-à-pair est plus 
difficile.

\paragraph{Cas d'un réseau point-à-point.} Dans le cas d'un réseau 
point-à-point, lorsque les deux pairs (\(n=2\)) sont connectés, alors \(p=1\). 
Les deux pairs forment, par nature, une composante géante, avec \(\alpha=0\) et 
un seul chemin. Ce qui implique nécessairement qu'il y a un point de 
défaillance unique. La démonstration résultante est donc moins intéressante. 
Mais cette défaillance se produit uniquement pour ces deux pairs et non pour le 
réseau entier~:~le reste des connexions continent de fonctionner.

\printbibliography

\end{document}